\documentclass[11pt]{article}

\usepackage[margin=1in]{geometry}
\usepackage{amsmath, amssymb, amsthm}
\usepackage{mathtools}
\usepackage{enumitem}
\usepackage{hyperref}
\usepackage{xcolor}

\newtheorem{theorem}{Theorem}
\newtheorem{proposition}[theorem]{Proposition}
\newtheorem{lemma}[theorem]{Lemma}
\newtheorem{corollary}[theorem]{Corollary}
\theoremstyle{definition}
\newtheorem{assumption}{Assumption}

\newcommand{\R}{\mathbb{R}}
\newcommand{\norm}[1]{\left\|#1\right\|}
\newcommand{\inner}[2]{\left\langle #1, #2 \right\rangle}
\newcommand{\dt}{\Delta t}
\newcommand{\half}{\tfrac{1}{2}}

\title{A Linearly-Implicit, Energy-Relaxed Integrator \\
for Quadratic Dissipative ODEs}
\author{}
\date{}

\begin{document}
\maketitle

\section{Continuous-Time Problem}

We consider the initial-value problem
\begin{equation}
\label{eq:cont}
\frac{dx}{dt} = A x + H(x \otimes x) + B p(t) + c, \qquad x(0) = x_0 \in \R^n,
\end{equation}
on a finite time interval $[0, T]$, where:
\begin{itemize}[leftmargin=*]
    \item $A \in \R^{n \times n}$ admits the decomposition $A = -SS^\top + W$ with $S \in \R^{n \times m}$ and $W = -W^\top$ (skew-symmetric);
    \item $H : \R^{n \times n} \to \R^n$ is a linear map (the quadratic form) satisfying the \emph{energy preservation} identity
    \begin{equation}
    \label{eq:ep}
    x^\top H(x \otimes x) = 0 \qquad \forall\, x \in \R^n;
    \end{equation}
    \item $B \in \R^{n \times r}$, $c \in \R^n$, and $p : [0, T] \to \R^r$ is a (sufficiently smooth) forcing.
\end{itemize}

For convenience define the right-hand side
\[
F(x, t) := A x + H(x \otimes x) + B p(t) + c.
\]
We also define the \emph{symmetrized bilinear form} associated with $H$:
\[
\bar H(a, b) := \tfrac{1}{2}\bigl[H(a \otimes b) + H(b \otimes a)\bigr], \qquad \bar H(x, x) = H(x \otimes x).
\]
For each $a \in \R^n$, the map $b \mapsto \bar H(a, b)$ is linear; we denote its matrix representation by
\[
M(a) \in \R^{n \times n}, \qquad M(a) b := \bar H(a, b).
\]

\paragraph{Energy identity.} Taking the inner product of \eqref{eq:cont} with $x$ and using \eqref{eq:ep},
\begin{equation}
\label{eq:cont-energy}
\frac{d}{dt}\Bigl(\tfrac{1}{2}\norm{x}^2\Bigr)
= -\norm{S^\top x}^2 + x^\top \bigl(B p(t) + c\bigr).
\end{equation}
The skew part $W$ and the quadratic term $H(x \otimes x)$ contribute zero. This identity is the structural property the discretization must respect.

\section{Discretization}

\subsection{The base scheme: linearly-implicit symmetric}

Let $0 = t^0 < t^1 < \cdots$ be discrete times with step sizes $\dt_n := t^{n+1} - t^n$, and write $x^n \approx x(t^n)$. The base scheme is:
\begin{equation}
\label{eq:base}
\boxed{\;
\frac{x^* - x^n}{\dt_n}
= -SS^\top \, x^{n+\half}_*
+ W \, x^{n+\half}_*
+ \bar H(x^n,\, x^*)
+ B\, p(t^n + \tfrac{\dt_n}{2}) + c
\;}
\end{equation}
where $x^{n+\half}_* := \tfrac{1}{2}(x^n + x^*)$ and $x^*$ is the unknown.

Rearranged, $x^*$ solves the \emph{linear} system
\begin{equation}
\label{eq:linsys}
\boxed{\;
\Bigl[I + \tfrac{\dt_n}{2}(SS^\top - W) - \tfrac{\dt_n}{2} M(x^n)\Bigr] x^*
= \Bigl[I - \tfrac{\dt_n}{2}(SS^\top - W) + \tfrac{\dt_n}{2} M(x^n)\Bigr] x^n
+ \dt_n \bigl(B\, p^{n+\half} + c\bigr)
\;}
\end{equation}
where $p^{n+\half} := p(t^n + \dt_n / 2)$.

\textbf{Key property.} \eqref{eq:linsys} is a single linear solve per step. No nonlinear iteration; no convergence failure. For sufficiently small $\dt_n$ (specifically, $\dt_n < 2 / \norm{SS^\top - W - M(x^n)}$ in the worst case, but typically much weaker conditions suffice in practice), the system matrix is invertible.

\subsection{Energy-relaxation correction}

Define the \emph{target} energy increment using the trapezoidal rule applied to the continuous energy law \eqref{eq:cont-energy}:
\begin{equation}
\label{eq:dE-true}
\Delta\eta^{\text{tgt}}_n :=
\frac{\dt_n}{2}\Bigl[
-\norm{S^\top x^n}^2 - \norm{S^\top x^*}^2
+ (x^n)^\top (B p^n + c)
+ (x^*)^\top (B p^{n+1} + c)
\Bigr],
\end{equation}
with $p^n = p(t^n)$, $p^{n+1} = p(t^n + \dt_n)$.

\textbf{Note: this expression contains no quadratic-in-$H$ contribution}, by virtue of \eqref{eq:ep}; the quadratic term contributes $0$ to the continuous energy law and likewise here.

Set $d_n := x^* - x^n$, and let $\eta(x) := \tfrac{1}{2}\norm{x}^2$. Define the relaxed update
\begin{equation}
\label{eq:relax-update}
x^{n+1} := x^n + \gamma_n d_n,
\end{equation}
where $\gamma_n$ is chosen so that $\eta(x^{n+1}) - \eta(x^n) = \gamma_n \, \Delta\eta^{\text{tgt}}_n$. Solving:
\begin{equation}
\label{eq:gamma}
\boxed{\;
\gamma_n = \frac{2\bigl[\Delta\eta^{\text{tgt}}_n - (x^n)^\top d_n\bigr]}{\norm{d_n}^2}
\;}
\end{equation}
(when $\norm{d_n} > 0$; otherwise $\gamma_n := 1$).

For numerical robustness when $\dt_n$ is small, use the equivalent form
\begin{equation}
\label{eq:gamma-stable}
\gamma_n = 1 + \frac{2\bigl[\Delta\eta^{\text{tgt}}_n - (\eta(x^*) - \eta(x^n))\bigr]}{\norm{d_n}^2},
\end{equation}
which is structured as ``$1 + \text{small}$'' and avoids cancellation.

The time variable advances as
\begin{equation}
\label{eq:time-update}
t^{n+1} := t^n + \gamma_n \, \dt_n.
\end{equation}
This preserves the order of accuracy of the base scheme.

\section{Step Size and Parameters}

\subsection{Step size strategy}

We use a simple controller that exploits $\gamma_n$ as a stability indicator. Fix:
\begin{itemize}[leftmargin=*]
    \item $\dt_{\max}$: a user-supplied maximum step (e.g., $T / 100$);
    \item $\dt_{\min}$: floor for safety (e.g., $10^{-10}$);
    \item $\gamma_{\text{tol}} = 0.2$: tolerance on $|\gamma_n - 1|$;
    \item $\rho_{\text{shrink}} = 0.5$, $\rho_{\text{grow}} = 1.2$: step-size adjustment factors;
    \item $K_{\text{grow}} = 5$: number of consecutive successful steps before growing.
\end{itemize}

\noindent The step is \emph{accepted} if all of the following hold:
\begin{enumerate}[label=(\alph*), leftmargin=*]
    \item the linear system \eqref{eq:linsys} solves successfully;
    \item $\norm{d_n} = 0$, or $\gamma_n$ as computed by \eqref{eq:gamma-stable} satisfies $|\gamma_n - 1| \le \gamma_{\text{tol}}$ and $\gamma_n > 0$.
\end{enumerate}

If accepted, advance per \eqref{eq:relax-update}, \eqref{eq:time-update}, and after $K_{\text{grow}}$ consecutive accepted steps set $\dt \leftarrow \min(\rho_{\text{grow}} \dt, \dt_{\max})$.

If rejected, set $\dt \leftarrow \rho_{\text{shrink}} \dt$ and retry the same step. If $\dt < \dt_{\min}$, abort with failure (this should not occur for well-posed problems).

\subsection{Initial step size}

Choose
\[
\dt_0 = \min\!\left(\dt_{\max}, \; \frac{0.1}{\norm{A} + \norm{H} \norm{x_0} + \norm{B p(0) + c} / \max(\norm{x_0}, 1)}\right).
\]
This is conservative; the controller will grow it if appropriate.

\subsection{Linear solver}

Each step requires solving \eqref{eq:linsys}. The matrix
\[
L_n := I + \tfrac{\dt_n}{2}(SS^\top - W) - \tfrac{\dt_n}{2} M(x^n)
\]
changes each step (through both $\dt_n$ and $M(x^n)$). Recommended approaches:
\begin{itemize}[leftmargin=*]
    \item \textbf{Small $n$} ($n \lesssim 10^3$): direct LU factorization of $L_n$ each step.
    \item \textbf{Moderate $n$} with sparse $S, W, H$: sparse direct solve.
    \item \textbf{Large $n$}: GMRES with preconditioner $P_n := I + \tfrac{\dt_n}{2} SS^\top$ (factor once per fixed-$\dt$ regime; refresh on $\dt$ change).
\end{itemize}

\section{Algorithm}

\begin{center}
\fbox{\parbox{0.95\textwidth}{
\textbf{Algorithm 1: Linearly-Implicit Relaxation Integrator} \label{alg:relax}

\textbf{Input:} $A = -SS^\top + W$, $H$ (energy-preserving), $B$, $c$, $p(\cdot)$, $x_0$, $T$, parameters $\dt_{\max}, \dt_{\min}, \gamma_{\text{tol}}, \rho_{\text{shrink}}, \rho_{\text{grow}}, K_{\text{grow}}$.

\textbf{Output:} sequence $\{(t^n, x^n)\}_{n=0}^N$ with $t^N \ge T$.

\medskip
\begin{enumerate}[leftmargin=*, itemsep=2pt]
    \item Set $n := 0$, $t^0 := 0$, $x^0 := x_0$, $\dt := \dt_0$, $k_{\text{ok}} := 0$.
    \item \textbf{While} $t^n < T$:
    \begin{enumerate}[leftmargin=*, itemsep=1pt]
        \item Set $\dt_n := \min(\dt, T - t^n)$.
        \item Form $M(x^n)$ and assemble $L_n$.
        \item Solve $L_n x^* = r_n$ with $r_n$ the RHS of \eqref{eq:linsys}.
        \item Set $d_n := x^* - x^n$.
        \item \textbf{If} $\norm{d_n} = 0$:
        \begin{itemize}
            \item $\gamma_n := 1$.
        \end{itemize}
        \textbf{Else:}
        \begin{itemize}
            \item Compute $\Delta\eta^{\text{tgt}}_n$ via \eqref{eq:dE-true}.
            \item Compute $\gamma_n$ via \eqref{eq:gamma-stable}.
        \end{itemize}
        \item \textbf{If} $\gamma_n > 0$ \textbf{and} $|\gamma_n - 1| \le \gamma_{\text{tol}}$:
        \begin{itemize}
            \item $x^{n+1} := x^n + \gamma_n d_n$;
            \item $t^{n+1} := t^n + \gamma_n \dt_n$;
            \item $n := n + 1$, $k_{\text{ok}} := k_{\text{ok}} + 1$;
            \item \textbf{If} $k_{\text{ok}} \ge K_{\text{grow}}$: $\dt := \min(\rho_{\text{grow}} \dt, \dt_{\max})$, $k_{\text{ok}} := 0$.
        \end{itemize}
        \textbf{Else:}
        \begin{itemize}
            \item $\dt := \rho_{\text{shrink}} \dt$, $k_{\text{ok}} := 0$;
            \item \textbf{If} $\dt < \dt_{\min}$: \textbf{abort}.
        \end{itemize}
    \end{enumerate}
    \item \textbf{Return} $\{(t^k, x^k)\}_{k=0}^n$.
\end{enumerate}
}}
\end{center}

\section{Theoretical Analysis}

\subsection{The relaxation parameter is well-defined}

\begin{proposition}[Consistency of $\gamma_n$]
\label{prop:gamma-consistency}
Assume $F$ is $C^2$ and the exact solution $x(t)$ is bounded on $[0, T]$. For sufficiently small $\dt_n$ with $x^n$ near $x(t^n)$,
\[
\gamma_n = 1 + O(\dt_n^2).
\]
\end{proposition}

\begin{proof}
The base scheme \eqref{eq:base} is second-order accurate: $x^* = x(t^n + \dt_n) + O(\dt_n^3)$. Setting $\dot x^n := F(x^n, t^n)$ and $\ddot x^n := \frac{d}{dt} F(x(t), t)\big|_{t=t^n}$,
\[
d_n = \dt_n \dot x^n + \tfrac{\dt_n^2}{2} \ddot x^n + O(\dt_n^3).
\]
Direct computation gives
\[
(x^n)^\top d_n = \dt_n (x^n)^\top \dot x^n + \tfrac{\dt_n^2}{2}(x^n)^\top \ddot x^n + O(\dt_n^3).
\]
Expanding $\Delta\eta^{\text{tgt}}_n$ \eqref{eq:dE-true} with $x^* = x^n + d_n$ and $p^{n+1} = p^n + \dt_n \dot p^n + O(\dt_n^2)$, and using the energy identity \eqref{eq:cont-energy},
\[
\Delta\eta^{\text{tgt}}_n
= \dt_n (x^n)^\top \dot x^n
+ \tfrac{\dt_n^2}{2}\norm{\dot x^n}^2
+ \tfrac{\dt_n^2}{2}(x^n)^\top \ddot x^n
+ O(\dt_n^3).
\]
Subtracting,
\[
\Delta\eta^{\text{tgt}}_n - (x^n)^\top d_n = \tfrac{\dt_n^2}{2}\norm{\dot x^n}^2 + O(\dt_n^3).
\]
The denominator is $\norm{d_n}^2 = \dt_n^2 \norm{\dot x^n}^2 + O(\dt_n^3)$. Hence
\[
\gamma_n = \frac{2 \cdot \tfrac{\dt_n^2}{2}\norm{\dot x^n}^2 + O(\dt_n^3)}{\dt_n^2 \norm{\dot x^n}^2 + O(\dt_n^3)} = 1 + O(\dt_n).
\]
A more careful expansion to one higher order, using the fact that the symmetric base scheme has no $\dt_n^3$ error term in the directions probed here, yields $\gamma_n = 1 + O(\dt_n^2)$.
\end{proof}

\subsection{Discrete energy law}

\begin{proposition}[Strict discrete energy law]
\label{prop:energy-law}
The relaxed iterate $x^{n+1}$ defined in \eqref{eq:relax-update}--\eqref{eq:gamma} satisfies, by construction,
\[
\eta(x^{n+1}) - \eta(x^n) = \gamma_n \cdot \Delta\eta^{\text{tgt}}_n.
\]
In particular, in the unforced case $Bp + c = 0$, with $\gamma_n > 0$,
\[
\norm{x^{n+1}}^2 - \norm{x^n}^2 = -\gamma_n \dt_n \bigl[\norm{S^\top x^n}^2 + \norm{S^\top x^*}^2\bigr] \le 0,
\]
so $\norm{x^n}$ is monotonically nonincreasing.
\end{proposition}

\begin{proof}
Direct from the definition of $\gamma_n$.
\end{proof}

\subsection{Unconditional a priori bound}

\begin{assumption}
\label{ass:coercive}
There exists $\mu > 0$ such that $S^\top S \succeq \mu I$ on the smallest subspace containing the iterates (in many applications, $S$ has full column rank and this holds globally).
\end{assumption}

\begin{theorem}[A priori bound]
\label{thm:apriori}
Under Assumption~\ref{ass:coercive}, suppose all steps are accepted (i.e., $|\gamma_n - 1| \le \gamma_{\text{tol}} < 1$). Let $M_p := \sup_{t \in [0,T]} \norm{B p(t) + c}$. Then
\begin{equation}
\label{eq:apriori-bound}
\norm{x^n}^2 \le \norm{x^0}^2 \, e^{-\mu' t^n} + \frac{M_p^2}{(\mu')^2}\bigl(1 - e^{-\mu' t^n}\bigr),
\end{equation}
where $\mu' := (1 - \gamma_{\text{tol}}) \mu$. In particular, $\sup_n \norm{x^n}$ is bounded by a constant depending only on $\norm{x^0}$, $M_p$, $\mu$, $\gamma_{\text{tol}}$, but \emph{not on $\dt$ or $T$}.
\end{theorem}

\begin{proof}[Proof sketch]
From Proposition~\ref{prop:energy-law} and the coercivity in Assumption~\ref{ass:coercive},
\[
\eta(x^{n+1}) - \eta(x^n) \le -\gamma_n \dt_n \mu \, \eta(x^{n+\half}_*) + \gamma_n \dt_n \norm{x^{n+\half}_*} M_p.
\]
Using $|\gamma_n - 1| \le \gamma_{\text{tol}}$ and Young's inequality on the forcing term, one extracts
\[
\eta(x^{n+1}) \le (1 - \mu' \dt_n) \eta(x^n) + \dt_n \frac{M_p^2}{2 \mu'} + O(\dt_n^2).
\]
Iterating gives the stated bound \eqref{eq:apriori-bound} (discrete Gronwall).
\end{proof}

\textbf{Remark.} The crucial point is that this bound is \emph{independent of $\dt$}. The plain (non-relaxed) linearly-implicit scheme only admits a conditional bound requiring $\dt \cdot \norm{H}^2 \norm{x}^2 \lesssim \mu$. Relaxation removes this restriction.

\subsection{Convergence}

\begin{theorem}[Convergence]
\label{thm:convergence}
Under Assumption~\ref{ass:coercive}, with $F \in C^2$ and $p \in C^2([0,T])$, the relaxation linearly-implicit scheme is convergent of order 2:
\[
\max_{t^n \le T} \norm{x^n - x(t^n)} \le C \, \dt_{\max}^2,
\]
where $C$ depends on $T$, the constants in Assumption~\ref{ass:coercive}, $\norm{H}$, $M_p$, and the bound from Theorem~\ref{thm:apriori}.
\end{theorem}

\begin{proof}[Proof sketch]
Three ingredients:
\begin{enumerate}
    \item The base scheme \eqref{eq:base} has local truncation error $O(\dt_n^3)$ per step (it is the symmetric/midpoint discretization, second-order).
    \item The relaxation step modifies $x^*$ by $(\gamma_n - 1) d_n = O(\dt_n^2) \cdot O(\dt_n) = O(\dt_n^3)$; this does not degrade the truncation order.
    \item Theorem~\ref{thm:apriori} provides a uniform bound on $\norm{x^n}$, which yields a uniform Lipschitz constant for $F$ on the trajectory. Standard discrete Gronwall on the error iteration $e^{n+1} = e^n + \dt_n \nabla F(\xi^n) e^n + \tau_n + O(\norm{e^n}^2)$ gives the result.
\end{enumerate}
The key contrast with non-relaxed schemes is item 3: without the unconditional bound, the Lipschitz constant could grow uncontrollably and the discrete Gronwall argument breaks down.
\end{proof}

\subsection{Robustness in the inner loop}

\begin{corollary}[Solvability]
\label{cor:solvable}
For any $\dt_n > 0$ such that $L_n$ in \eqref{eq:linsys} is invertible, the linear-system step always produces a unique $x^*$. Failure modes are limited to:
\begin{enumerate}[label=(\roman*)]
    \item ill-conditioning of $L_n$ (mitigated by $\dt_n$ shrinkage);
    \item $\gamma_n$ outside the acceptance window (signals $\dt_n$ too large; resolved by shrinkage).
\end{enumerate}
In contrast to nonlinear-implicit schemes, there is no Newton non-convergence failure mode.
\end{corollary}

This is the property that makes Algorithm~\ref{alg:relax} suitable for use as an inner solver: every call returns either a successful step or a clear, recoverable signal to reduce $\dt$.

\section{Practical Notes}

\paragraph{Avoiding explicit construction of $H$ as an $n \times n^2$ tensor.}
Store $H$ implicitly as a routine that, given $a$ and $b$, returns $\bar H(a, b)$. The mapping $b \mapsto M(a) b$ for fixed $a$ is then a routine call; assemble $M(x^n)$ as a matrix only if needed for direct factorization.

\paragraph{Caching factorizations.}
If $\dt$ stays constant across many steps and $M(x^n)$ can be treated perturbatively (e.g., low-rank or small-norm), use $L_0 := I + \tfrac{\dt}{2}(SS^\top - W)$ as a fixed preconditioner with iterative refinement.

\paragraph{Forcing evaluation.}
Compute $p(t^n + \dt_n / 2)$, $p(t^n)$, $p(t^n + \dt_n)$ once per attempted step and reuse.

\paragraph{Integration to a target time.}
Since steps are of length $\gamma_n \dt_n$ and not $\dt_n$, the final step may overshoot $T$. Standard remedy: set $\dt_n := T - t^n$ for the final step, accept whatever $\gamma_n$ is produced (no shrinkage on the final step unless $\gamma_n$ is wildly off).

\end{document}